\chapter*{前言}
\addcontentsline{toc}{chapter}{前言}
\label{ch:preface}

这本书试图回答一个在 PCSEL 学习过程中反复出现、却经常得不到系统回答的问题：如果一个读者只有基础电动力学、基础光学和基础半导体激光器背景，那么他应该沿着怎样一条尽量连续、尽量少跳步的路径，才能真正建立起对 photonic-crystal surface-emitting laser（PCSEL，光子晶体面发射激光器）的完整理解？这里所谓“完整理解”，并不是指记住若干术语，也不是停留在几张典型结构图和若干仿真截图上，而是能够把物理图像、数学模型、数值方法和器件实现联系起来，知道每一个模型在回答什么问题，也知道每一个模型不能替代什么问题。

PCSEL 之所以值得单独写成一本书，是因为它天然跨越了多个通常被分开教授的知识模块。把它仅仅视为一个光子晶体共振器，会忽略它作为半导体激光器的增益、损耗、阈值、电注入和热效应；把它仅仅视为一个半导体激光器，又会错过二维周期开放电磁系统、Bloch 模、Gamma 点、辐射耦合和远场工程这些核心特征。真正的 PCSEL 理解，必须同时包含三个层面：它是一个周期开放电磁系统；它是一个具有增益与损耗竞争的有源光学腔；它还是一个依赖外延层、载流子输运、电流扩展和热分布的真实半导体器件。只抓住其中任意一个层面，结论都会变得片面，甚至在设计上导向错误。

因此，本书的组织原则不是“把所有相关名词列一遍”，而是尽量把它们压缩成一条可连续阅读的主线。前几章从 Maxwell 方程、周期介质和 Bloch 理论出发，先建立读者理解 PCSEL 所需的电磁学骨架；随后讨论 Gamma 点、带边、辐射耦合和面内反馈，把“为什么能面发射、为什么可能大面积单模、为什么不能把被动高 $Q$ 直接等价成低阈值”这些最容易被误解的问题讲清楚；再往后进入耦合波理论、平面波展开法、RCWA、FDTD 与 FEM，使读者能够从不同保真度层次理解和计算 PCSEL；后半部分则把重点转向半导体外延层、量子阱增益、电注入、漂移--扩散、热效应和多物理场耦合，说明为什么真实器件设计不能停留在“冷腔 + 均匀增益”的理想模型上。

我在写作中刻意坚持几个原则。第一，任何一个重要结论都尽量给出来源：要么给出完整推导，要么清楚说明证明思路和省略步骤，而不是只报结论。第二，任何一个常用近似都尽量说明适用条件，例如 slab 中的近似 TE/TM 分类、effective index 近似、均匀增益模型、有限傅里叶截断、冷腔阈值代理等都不是无条件成立的。第三，任何一个仿真结论都必须说明它来自哪一类求解问题，回答的是单胞、无限周期、有限腔、被动开放系统还是有源器件问题。第四，我反复强调阈值增益与阈值电流、被动共振与实际激射模、能带候选模与真实工作模之间的区别，因为这正是初学者最容易混淆、也是研究中最容易导致错误判断的几个节点。

这本书并不追求把每一种商用软件的操作界面讲成“按钮教程”。软件会更新，界面会变化，真正稳定的是控制方程、边界条件、离散思想、收敛标准以及不同求解器之间的数据接口。因此，书中关于数值方法的部分始终坚持 solver-agnostic 的写法：首先解释一个方法在数学上解的是什么；其次说明几何、材料和边界条件如何进入模型；再说明输出量怎样与 PCSEL 的设计问题发生联系。在这个框架下，我也适度给出 Lumerical 和 COMSOL 一类工具的典型使用分工，目的是帮助读者建立实际工作流，而不是把理解替换成软件熟练度。

如果你是第一次接触 PCSEL，最稳妥的阅读方式是顺着本书章节次序向前推进，不要一开始就跳到器件级阈值或多物理场仿真章节。因为那些章节虽然看起来更接近工程结果，但它们的输入量，例如候选模、约束因子、材料增益、损耗分布、热反馈灵敏度，本质上都建立在前面的电磁与模式分析之上。如果基础骨架没有搭起来，后面的器件结果只会变成一堆参数和软件设置。如果你已经具备一定背景，则可把本书当作一份交叉索引式参考：围绕某个具体问题，例如“Gamma 点模的辐射损耗如何调控”“为什么高 $Q$ 模不一定是最低阈值模”“电流扩展怎样影响模式竞争”，沿着书中的交叉引用回溯相关章节。

我也必须坦率说明本书的边界。它不是某一套实验数据的总结，也不是某一类具体材料体系的工艺手册。书中不会伪造实验结果，不会虚构数据库参数，也不会把未经验证的近似公式包装成严格理论。凡是需要材料数据库、精确工艺窗口或特定软件版本支持的地方，本书都只提供建模框架、参数类型、数量级意识和验证路径，而不冒充已经给出唯一正确答案。对于 PCSEL 这样跨越电磁学、半导体器件物理和数值计算的主题，这种克制是必要的，因为只有明确知道哪些内容是严格的、哪些内容是近似的、哪些内容仍然需要校准和实验闭环，书稿才可能真正成为可靠的研究起点。

如果这本书能够帮助读者从“会看一些图和公式”进一步走到“能够独立判断一个 PCSEL 模型是否合理、一个仿真流程是否闭合、一个器件结论是否越界”，那么它就达到了写作目的。希望读者在阅读过程中始终记住一件事：PCSEL 不是抽象的光场图样，也不是孤立的能带曲线，而是一个把周期电磁结构、有源增益介质和真实半导体器件约束同时绑定在一起的问题。只有在这三个层面上同时保持清醒，后续的设计、仿真与研究判断才会真正可靠。

\begin{flushright}
Feiyang Wu\\
2026 年 3 月
\end{flushright}
